\documentclass[12pt, a4paper]{article}

\usepackage{amsmath, amsthm, amssymb}
\usepackage{titling}
\usepackage{geometry}
\usepackage{setspace}
\usepackage{titlesec}
\usepackage{hyperref}
\usepackage{booktabs}
\usepackage{array}
\usepackage[utf8]{inputenc}
\usepackage{parskip}

\geometry{margin=1.25in}
\onehalfspacing

\hypersetup{
  colorlinks=true,
  linkcolor=black,
  citecolor=black,
  urlcolor=black
}

\theoremstyle{plain}
\newtheorem{proposition}{Proposition}[section]
\newtheorem{corollary}[proposition]{Corollary}
\newtheorem{lemma}[proposition]{Lemma}
\newtheorem{hypothesis}[proposition]{Hypothesis}

\theoremstyle{definition}
\newtheorem{definition}[proposition]{Definition}
\newtheorem{assumption}[proposition]{Assumption}

\theoremstyle{remark}
\newtheorem{remark}[proposition]{Remark}

\titleformat{\section}{\large\bfseries}{\thesection.}{0.75em}{}
\titleformat{\subsection}{\normalsize\bfseries}{\thesubsection.}{0.75em}{}

\title{\Large\textbf{Compression After Expansion:\\
On Skill, Misattribution, and the Ethics of Visible Structure}}
\author{Flyxion\\
\small Independent Researcher}
\date{2026}

\setlength{\droptitle}{-3cm}
\posttitle{\par\end{center}\vspace{0.3cm}}
\predate{\begin{center}\normalsize}
\postdate{\par\end{center}\vspace{0.2cm}}

\begin{document}

\maketitle

\begin{abstract}
This essay develops a qualified model of a recurring pattern in human skill, technical production, and social representation: apparent simplicity in a finished artifact---a mathematical proof, a constructed object, a piece of software, a fluent performance---is often the residue of prior exploration, iteration, and structural accumulation that the finished form no longer displays. We formalize a narrow version of this claim under explicit assumptions about a constraint-defined configuration space, rather than asserting it as a universal structural law, and we distinguish the resulting probabilistic result from claims about relative difficulty, which require further argument. We then develop Aspect Relegation Theory, a cognitive model in which a substantial class of expert automatic performance is understood as the residue of previously deliberate procedures that became increasingly autonomous through practice; we present this as one well-supported route to automaticity among others, not as an identity claim about all fast cognition. A candidate neurophysiological correlate---transient hypofrontality in flow states---is discussed as a contested hypothesis rather than an established signature, alongside evidence that complicates a simple prefrontal-withdrawal story. These two frameworks motivate an account of observer misattribution, of the political economy of delegated constraint management, and of the ethics of compressed representation, organized around a narrower and more defensible thesis than the one with which earlier drafts of this argument began: that constructing a reliable representation of a domain generally requires prior exploration of that domain by someone, and that observers who encounter only the compressed result routinely miscalibrate the labor required to produce it.
\end{abstract}

\tableofcontents

\newpage

\section{Introduction: The Illusion of Effortless Form}
\label{sec:intro}

There is a persistent pattern in the way human beings perceive completed artifacts. A finished building, a published proof, a released piece of software, a polished essay---each presents itself as a coherent, apparently simple object whose internal structure invites little reflection on the process by which it came to be. Observers encountering such objects tend, with some consistency, to draw inferences about the difficulty of producing them that track the apparent simplicity of the result rather than the generative complexity of the process. This inference is often badly calibrated, and understanding why is the primary aim of this essay. It is not, however, a universal law: some finished objects really were easy to make, and surface simplicity sometimes does correlate with low production cost. The claim developed here is about a common and consequential bias, not an exceptionless regularity.

The illusion has several reinforcing sources. Compression is a constitutive feature of representation: the point of producing a final form is often precisely to eliminate from view the intermediate steps, failed attempts, and structural discoveries that preceded it. This is not a contingent feature of bad communication but a functional requirement of efficient representation. A proof that displayed every exploratory conjecture, failed lemma, and notational revision through which it was produced would be unreadable. A building that preserved every scaffolding stage would be uninhabitable. Compression is, in this sense, often necessary. But it can also be epistemically misleading when the compressed form is mistaken for the primary form, and when the prior expansion is forgotten or invisible.

A second reinforcing source is cognitive. Human beings develop many skills through a process that renders previously effortful operations increasingly automatic, relocating them beneath the threshold of conscious attention. Expert performance can therefore appear effortless because, for the expert, much of it is---but this effortlessness is frequently the product of accumulated practice, not evidence of an absence of underlying complexity. The experienced builder who walks through a derelict structure and immediately perceives a tractable sequence of operations has a learned decomposition that converts visible disorder into a sequence of workable steps; the novice who sees only disorder is not mistaken about the building's current state, only about how it might be reached from here.

A third source is institutional and economic. Modern media and commercial systems are organized substantially around the representation of outcomes rather than processes. Startup narratives, viral content, minimalist design aesthetics, and compressed pedagogical formats often function by presenting relegated forms while suppressing the construction histories that produced them. This is not always deliberate deception, but it can reliably produce distortions in how difficulty, effort, and value are perceived and distributed.

The core claim of this essay may be stated directly, with the qualification the rest of the essay will maintain: \textit{much of what appears simple is the result of prior, invisible complexity that has been compressed}, and this compression presupposes some prior exploration, is achieved through stabilization, and is fragile in ways the compressed form does not reveal on its own. This is a claim about a common structural pattern with real consequences for labor, pedagogy, and the ethics of representation---not a claim that every simple-looking thing has a hidden history, or that all forms of ease are equally deceptive.

\section{Expansion, Compression, and a Qualified Formal Model}
\label{sec:formal}

\subsection{Preliminaries}
\label{subsec:preliminaries}

We introduce a minimal formal vocabulary sufficient to state a narrow, conditional version of the structural asymmetry between construction and compression. The vocabulary is offered as a model under explicit assumptions, not as a description that automatically holds of every constrained system.

\begin{definition}
Let a \emph{system} $\mathcal{S}$ be a configuration of components subject to a constraint set $\mathcal{C}$. A configuration $s$ is \emph{coherent} if and only if $s \models \mathcal{C}$.
\end{definition}

Let $\Omega$ denote the space of possible configurations of $\mathcal{S}$, and let
\[
\Omega_{\mathrm{valid}} = \{ s \in \Omega \mid s \models \mathcal{C} \}, \qquad \Omega_{\mathrm{invalid}} = \Omega \setminus \Omega_{\mathrm{valid}}.
\]

Cardinality is the wrong instrument for comparing these sets whenever $\Omega$ is infinite, since a low-measure subset and its ambient space can share the same cardinality, and the raw statement $|\Omega_{\mathrm{valid}}| \ll |\Omega|$ is not a structural truth about constrained systems in general: a constraint set can be weak, redundant, or highly correlated, and can leave a large valid region even when numerous. We therefore fix a measure $\mu$ on $\Omega$ appropriate to the system under discussion---counting measure for finite discrete configuration spaces, a natural probability or volume measure otherwise---and state the assumption this essay actually needs explicitly.

\begin{assumption}[Sparse validity]
\label{ass:sparse}
For the systems considered in this essay,
\[
\frac{\mu(\Omega_{\mathrm{valid}})}{\mu(\Omega)} \ll 1.
\]
\end{assumption}

This is not asserted as a law governing all constrained systems; it is a condition satisfied by many of the domains the essay discusses---working software, load-bearing structures, valid proofs---and the results that follow are conditioned on it holding. Where it fails, the arguments of this section do not apply, and that is a substantive limitation rather than a technicality.

\subsection{Expansion as Exploration}
\label{subsec:expansion}

\begin{definition}
\emph{Expansion} is a process $E$ acting on an initial configuration $s_0 \in \Omega$, generating a sequence $E(s_0) = \{s_1, s_2, \dots, s_n\}$ where each $s_i$ is obtained through transformation, trial, or variation.
\end{definition}

Expansion is characterized by traversal of configuration space, evaluation of constraint satisfaction, discovery of failure modes, and articulation of dependencies. It is therefore not merely generative but epistemic: expansion reveals structure in $\Omega_{\mathrm{valid}}$ through interaction with $\mathcal{C}$. This is the process through which invariants are discovered, dependency structures are mapped, and failure modes are identified and resolved. Expansion need not be performed by the agent who later benefits from it; it can be carried out by another person, an institution, a culture, an evolutionary process, or an earlier generation of a system, and its results transmitted rather than rediscovered.

\subsection{Compression as Projection}
\label{subsec:compression}

A bare mapping $C: \Omega_{\mathrm{valid}} \to \Sigma$ is too weak to capture what is meant by compression, since a constant function satisfies it while preserving nothing, and ``lower-dimensional'' is not a defined property of an arbitrary representational space $\Sigma$. We instead require a decoder and a preservation condition.

\begin{definition}
\emph{Compression} is a triple $(C, D, \sim)$ consisting of an encoder $C: \Omega_{\mathrm{valid}} \to \Sigma$, a decoder $D: \Sigma \to \Omega$, and a task-relevant equivalence relation $\sim$ on $\Omega_{\mathrm{valid}}$, such that for all $s \in \Omega_{\mathrm{valid}}$,
\[
D(C(s)) \sim s,
\]
possibly subject to a bounded distortion condition rather than exact equivalence.
\end{definition}

Compression under this definition eliminates redundancy while preserving task-relevant structure, encoding invariants discovered during expansion and suppressing intermediate derivations that are not required to reconstruct behavior up to $\sim$. Typical instances include refactored code, minimal mathematical proofs, concise explanatory prose, and completed architectural forms. In the cases this essay considers, compression chiefly represents structure discovered or acquired elsewhere, although the act of compressing something can itself expose further invariants---dimensionality reduction and model selection are genuine forms of discovery, not merely encoding, and nothing here denies that.

\subsection{The Dependency of Compression on Prior Exploration}
\label{subsec:dependency}

A universal claim that no valid compression can be constructed without expansion does not hold: a mapping might be inherited from another agent, derived analytically from already-known constraints, or discovered by accident without systematic traversal. The defensible version of the claim is epistemic and conditional.

\begin{proposition}[Conditional dependency of compression]
\label{prop:compression-requires-expansion}
In an initially unknown task domain, constructing a reliable representation $(C, D, \sim)$ generally requires information about which variations preserve $\sim$-relevant validity. Absent an external source of this information---prior search by the agent, observation, instruction, or an inherited body of structural knowledge from elsewhere---such information must be acquired through exploration of $\Omega$ near the boundary of $\Omega_{\mathrm{valid}}$.
\end{proposition}

\begin{proof}[Sketch]
A compression that removes components while preserving $\sim$ requires knowing which components are invariant under $\sim$ and which are incidental. Where this knowledge is not supplied externally, it can only be established by observing which variations preserve validity and which do not, which is exploration of the boundary between $\Omega_{\mathrm{valid}}$ and $\Omega_{\mathrm{invalid}}$ in the sense of Section~\ref{subsec:expansion}.
\end{proof}

The practical corollary this essay relies on is narrower than the earlier, stronger claim: one cannot, without access to some prior exploration---one's own or someone else's, transmitted through instruction, culture, or inherited design---produce a reliable compressed representation of a domain one does not yet understand. This applies to mathematical exposition, software architecture, physical construction, and written argument, though the exploration in question need not have been performed personally or recently by the person producing the final form.

\subsection{Compression Without Relegation: A Limiting Case}
\label{subsec:compression-without-relegation}

Proposition~\ref{prop:compression-requires-expansion} already licenses the possibility that the expansion behind a compression was performed by an agent other than the one presenting it. It is worth pushing this to its limit, because the limiting case clarifies a distinction the rest of this essay depends on: between an artifact's observable fluency and the developmental process, if any, that an agent using that artifact has itself undergone.

Consider an agent whose entire competence in a domain consists of an inherited representation $(C, D, \sim)$, where the agent performed none of the exploration in Section~\ref{subsec:expansion} that produced it. Large language models are the clean paradigm case. A model's fluent, low-latency completions are a compression over regularities discovered through the effortful, iterative, error-correcting production of the human text on which it was trained---genuine expansion, performed by many agents over a long history, near the boundary between $\Omega_{\mathrm{valid}}$ and $\Omega_{\mathrm{invalid}}$ for countless task domains. The model's forward pass exhibits the surface signature of relegated performance---speed, fluency, the absence of visible deliberation---but nothing in that forward pass corresponds to an agent's own trajectory from deliberate monitoring to earned stabilization, because no such trajectory occurred for the system producing the output. We will call this \emph{inherited compression}: a representation transmitted to an agent that performed none of the exploration behind it, as opposed to relegation, where the agent that now executes automatically is the same agent that once explored deliberately.

This distinction reframes a question that recurs whenever dual-process language is applied to AI systems: whether a model has ``System 1'' processing without a well-developed ``System 2.'' Posed that way, the question asks only about behavioral signature---is the output fast and non-deliberative---and answers built on that framing tend to list what the system lacks (emotional grounding, contextual understanding) without engaging the structural question underneath. Section~\ref{subsec:art-formal}'s model supplies that question directly: not whether behavior looks automatic, but whose trajectory produced its stabilization. Inherited compression can reproduce the behavioral signature of relegation---speed, fluency, apparent effortlessness---without instantiating the process the signature is usually taken to indicate. Treating fast, fluent output as evidence of relegated System 1 processing therefore risks a category error: it mistakes a property of the artifact's presentation for a property of the process that generated it, precisely the conflation Section~\ref{sec:invisibility} diagnoses in other contexts.

Inherited compression is not the only mechanism that can produce automaticity without the agent's own deliberate exploration. A shorter, contrasting case is instructive. When an advertiser transforms a brand into characters and repeats the pairing until consumers respond to the brand automatically, the consumer does possess the attentional machinery relegation requires---an active attention set, a capacity for stabilization. But the stabilization was not produced by the consumer's own deliberate practice at solving a task of theirs; it was produced by an external actor engineering the consumer's exposure history toward a fixed association, for that actor's objective rather than the consumer's own. We will call this \emph{conditioning}: stabilization that is genuinely the agent's own, but whose triggering exposure was authored by someone else. Section~\ref{subsec:type1type2} returns to this contrast with the full apparatus of Section~\ref{subsec:art-formal} once that apparatus has been introduced.

Together, the two cases sharpen the diagnostic question this essay's model is built to ask. It is not enough to observe that behavior is fast, fluent, or free of visible deliberation. The relevant questions are whose trajectory produced the underlying stabilization, by what mechanism it was produced, and for whose objective---the agent's own task, in the case of genuine relegation; another agent's prior task, transmitted rather than repeated, in the case of inherited compression; or a third party's objective, imposed on the agent's own stabilization machinery, in the case of conditioning.

\subsection{A Probabilistic Asymmetry Between Construction and Disruption}
\label{subsec:asymmetry}

An earlier formulation of this section defined a pseudo-entropy $H(s) \propto \log|\Omega_{\mathrm{reachable}}(s)|$ and claimed that the number of entropy-increasing operations exceeds the number of entropy-decreasing operations under any fixed constraint set. Neither claim survives scrutiny: $\Omega_{\mathrm{reachable}}(s)$ depends on an arbitrary choice of operations and horizon, and the counting claim is false as stated, since an operation repertoire consisting entirely of repair moves, or one in which every destructive operation is paired with a constructive inverse, is a live counterexample. What the argument of this essay actually needs is a probabilistic statement about perturbations, not a claim about the cardinality of an operation set.

\begin{proposition}[Perturbation asymmetry under sparse validity]
\label{prop:asymmetry}
Suppose Assumption~\ref{ass:sparse} holds, and let perturbations of a valid state be drawn from a distribution over $\Omega$ that is not concentrated on validity-preserving transformations. Then a randomly sampled perturbation of a valid configuration $s \in \Omega_{\mathrm{valid}}$ is more likely to leave $\Omega_{\mathrm{valid}}$ than to remain within it.
\end{proposition}

This is a genuine probabilistic consequence of sparse validity together with a stated assumption about the perturbation distribution; it is not, by itself, a claim that evaluation is easier than construction, and the two should not be run together. Detecting that a constraint has failed can itself be expensive---requiring measurement, proof, or diagnosis that is far from trivial in domains such as distributed systems or formal verification. Breaking a system, detecting that it has broken, explaining why, and repairing it are four distinct operations with different cost profiles, and collapsing them into an undifferentiated category of ``destructive and evaluative processes'' overstates what follows from Proposition~\ref{prop:asymmetry}.

What the proposition does support is narrower and still useful: under sparse validity and a perturbation distribution not aimed at preserving coherence, maintaining a system in $\Omega_{\mathrm{valid}}$ requires selecting from a comparatively small set of validity-preserving operations, while an unstructured perturbation is likely to exit that set. We will refer to $\Omega_{\mathrm{valid}}$ as a \emph{narrow admissible region} within $\Omega$ rather than describing it geometrically as a low-dimensional manifold, since a constraint-defined subset can be discrete, disconnected, fractal, combinatorial, or full-dimensional but small in measure; the geometric language of manifolds is not established by anything argued here.

\subsection{Fragility of Compressed Representations}
\label{subsec:fragility}

Because compression under Definition~\ref{subsec:compression} encodes a system's behavior only up to the equivalence relation $\sim$, compressed representations are fragile in a specific sense: perturbations that move a configuration outside the region the encoder was calibrated against can produce invalid decodings; omitted intermediate structure can prevent reconstruction of the generative pathway; and interpretability of the compressed form depends on the receiver's prior familiarity with the domain from which it was extracted.

This fragility grounds a further point about evaluative ease that recurs throughout this essay. Agents whose primary interaction with a system is evaluative---critics, reviewers, inspectors, users---typically interact with a system already inside $\Omega_{\mathrm{valid}}$, and their task is comparatively local: identify some perturbation that would move the system into $\Omega_{\mathrm{invalid}}$, or notice that it already has. Under Assumption~\ref{ass:sparse} and Proposition~\ref{prop:asymmetry}, this local task samples from a different, and typically larger, region of possibility space than the task of constructing or maintaining the system in the first place. The apparent ease of finding a flaw is therefore not, by itself, evidence that avoiding that flaw during construction was equally easy; the two tasks are not drawing from the same distribution of operations, though how much easier evaluation is remains an empirical question that varies by domain and is not settled by the formal result alone.

\begin{quote}
\itshape Under sparse validity, exploration is required to locate coherent structure, and compression encodes what exploration has found. Constructing or maintaining a valid configuration is movement into and continued occupancy of a narrow admissible region; a validity-breaking perturbation, by contrast, need only exit that region by any available route. This asymmetry is conditional on the assumptions stated above, not a universal law, but where those assumptions hold it helps explain why compressed artifacts, encountered without their generative histories, are so often misjudged.
\end{quote}

\section{Skill as Internalized Structure}
\label{sec:relegation}

\subsection{The Standard Account and Its Inadequacy}

The standard account of expertise treats skill as accumulated knowledge and trained performance. This is not wrong, but it is underspecified in a way that limits its ability to explain some of the most striking features of expert cognition: the disappearance of felt effort, the compression of explanation, and the gap between expert and novice perception of the same situation. A more adequate account requires attention to what happens to the structure of a task as it is learned, not merely to the content of what is known.

\subsection{Aspect Relegation Theory: A Revised Formal Model}
\label{subsec:art-formal}

An earlier version of this model set $|A(t)| \leq k$ for a bounded active attention set while also setting $A(0) = T$, the full task; these are jointly inconsistent unless $|T| \leq k$, and for a novice this is typically false---the novice's characteristic difficulty is precisely that the task has more aspects than can be simultaneously attended to. The revised model separates three things the earlier version conflated: the full task structure, what the agent has come to represent about it, and what the agent is currently attending to.

\begin{definition}
Let a \emph{task} $T = \{a_1, \dots, a_n\}$ be a finite set of aspects, each corresponding to a constraint, sub-operation, or dependency relevant to successful execution. Let $K(t) \subseteq T$ be the \emph{known aspects}: those the agent has represented at time $t$, whether or not currently monitored. Let the \emph{active attention set} $A(t) \subseteq K(t)$ be the aspects under conscious monitoring at time $t$, subject to a bounded capacity $|A(t)| \leq k$.
\end{definition}

\begin{definition}
An aspect $a_i$ is \emph{stabilized to level $1-\varepsilon$ under policy $\pi$ and environment class $\mathcal{E}$} if
\[
\Pr(\text{constraint satisfied} \mid a_i, \pi, \mathcal{E}) \geq 1 - \varepsilon.
\]
\end{definition}

Stabilization is thus probabilistic and context-relative rather than a binary property. This is what makes it possible for a stabilized aspect to fail outside the environment class in which it was learned, and it is why expert automaticity transfers unevenly across contexts: a routine reliable under $\mathcal{E}$ carries no guarantee under $\mathcal{E}'$.

\begin{definition}
The \emph{attention update rule} is
\[
A(t+1) = \bigl(A(t) \setminus R_t\bigr) \cup N_t \cup Q_t,
\]
where $R_t \subseteq A(t)$ is the set of aspects that have met their stabilization threshold and are relegated out of active attention; $N_t \subseteq K(t) \setminus A(t)$ is a set of newly attended higher-order aspects that become tractable once capacity is freed by $R_t$; and $Q_t$ is a set of aspects---possibly previously relegated---reactivated by uncertainty, novelty, unfamiliar context, or detected error.
\end{definition}

This rule replaces the earlier purely subtractive operator. Learning is not monotonic in this model: as lower-level operations stabilize, new higher-order aspects can enter $A(t)$ via $N_t$; previously stabilized routines can be recruited back into active attention via $Q_t$ when they fail outside their training distribution or when the agent detects an anomaly. We correspondingly drop the earlier limit claim $A(t) \to \emptyset$. Skilled performance typically retains nonzero attention at some level---goal selection, anomaly monitoring, strategic coordination, environmental feedback---and is better modeled as a change in \emph{what} is monitored than as the disappearance of monitoring altogether.

\subsection{Automaticity as a Class of Relegation Outcomes, Not a Universal Identity}
\label{subsec:type1type2}

The dual-process literature distinguishes \emph{Type 1} processing---rapid, autonomous, low-working-memory-load---from \emph{Type 2} processing, which supports hypothetical reasoning and loads working memory \cite{evans2013}. We use ``Type 1'' and ``Type 2'' rather than the popular ``System 1'' and ``System 2'' labels except when referring to that popular framing directly. This is more than a terminological nicety. Bargh's review of the automaticity literature found no single process that satisfies all of the criteria---unintentional, unconscious, efficient, uncontrollable---popularly bundled into ``automatic,'' undermining any picture of automaticity as one unified, box-checkable state that a system either has or lacks \cite{bargh1992}. Kahneman himself was explicit that System 1 and System 2 name types of processing rather than locations or literal systems, that ``there is no part of the brain that either of the systems would call home,'' and that the two are typically interleaved rather than strictly sequential \cite{kahneman2011}. Popular treatments routinely lose this qualification, treating the systems as anatomically real and asking whether a given process ``has'' one or the other as though that were a binary fact about the process rather than a claim about its history.

An earlier version of this essay claimed that ``System 1 is not a different kind of thinking; it is System 2 at a later stage of its own development''---an identity claim that the evidence does not support. Three distinct claims need to be separated.

\begin{enumerate}
\item[(i)] Controlled procedures can, through repetition and correction, become increasingly autonomous.
\item[(ii)] All expert intuition has controlled, deliberative origins.
\item[(iii)] All Type 1 cognition is compiled Type 2 cognition.
\end{enumerate}

Claim (i) is well supported: it is the content of the revised relegation model above, and it accounts for a large and important class of expert automaticity. Claim (ii) is doubtful, and claim (iii) is false. Automaticity research finds that automatic higher-cognitive processes do not arise exclusively from a history of conscious skill acquisition; some automatic processes have evolved substrates or are established through early childhood learning mechanisms that were never consciously decomposed by the individual who now exhibits them \cite{bargh2012}. Motor-learning research on explicit and implicit training finds that both routes can lead toward automatic, dual-task-robust performance, with the balance of evidence in some paradigms actually favoring implicit acquisition---routines that were never represented explicitly enough to be called relegated System 2 procedures in the first place \cite{kal2018}.

\begin{proposition}[Relegation as one route to Type 1 automaticity]
\label{prop:s1-as-s2}
A substantial class of expert Type 1 performances consists of procedures that once required controlled attention---were represented in $A(t)$ and governed by explicit monitoring---and became increasingly autonomous through the practice, chunking, and probabilistic stabilization described in Section~\ref{subsec:art-formal}. This class does not exhaust Type 1 cognition, which also includes innate, perceptual, affective, associative, and implicitly acquired processes with no comparable history of explicit control.
\end{proposition}

\begin{corollary}
Where an expert Type 1 performance does belong to the relegation class, what appears as immediate intuition is, for that performance specifically, the residue of prior deliberation. The corollary does not generalize to Type 1 cognition as a category.
\end{corollary}

\subsubsection*{Three mechanisms, restated formally}

Section~\ref{subsec:compression-without-relegation} distinguished inherited compression and conditioning from genuine relegation informally, before the apparatus of Section~\ref{subsec:art-formal} was available. With $K(t)$, $A(t)$, $R_t$, and the stabilization condition $\Pr(\text{constraint satisfied} \mid a_i, \pi, \mathcal{E}) \geq 1-\varepsilon$ now in hand, the three mechanisms can be stated on common ground.

\emph{Genuine relegation} holds when the stabilization probability is estimated over the agent's own trajectory through $K(t) \to A(t) \to R_t$: the same agent that once represented $a_i$ under active monitoring is the one whose repeated attempts, under its own policy $\pi$, drove the constraint-satisfaction probability above threshold. \emph{Inherited compression} holds when a receiving agent is given $(C, D, \sim)$ directly, without ever instantiating a trajectory through $\mathcal{E}$ of its own; there is no $A(t)$ from which anything was relegated, only an artifact whose stabilization---if the term even applies---was earned by a different agent entirely. \emph{Conditioning} holds when the stabilization probability is genuinely estimated over the receiving agent's own responses, but the policy $\pi$ and the distribution of exposures constituting $\mathcal{E}$ were selected by an external actor for that actor's objective rather than emerging from the agent's own deliberate attempts to solve a task. All three can produce output that is fast, fluent, and free of visible deliberation; only the first is relegation in the sense this essay uses the term, and the differences among the three concern reliability, transfer, and whose interests the resulting automaticity actually serves.

\subsection{Neural Correlates: A Contested Hypothesis, Not an Established Signature}
\label{subsec:hypofrontality}

If a substantial class of skilled automatic execution really does correspond to withdrawal of active monitoring from stabilized aspects, and if the prefrontal cortex is centrally implicated in that monitoring function, then some forms of flow-like performance might be expected to show reduced prefrontal engagement relative to the same task under full deliberative load. This is essentially the \emph{transient hypofrontality hypothesis}, proposed by Dietrich as a mechanism for the flow state and for altered states of consciousness more generally \cite{dietrich2003,dietrich2004}. It is important to be precise about its status: Dietrich's own papers present it as a hypothesis, not as an established finding, and later experimental work complicates the simple prefrontal-withdrawal story considerably.

\begin{hypothesis}[Qualified hypofrontality hypothesis]
\label{hyp:hypofrontality}
If flow involves reduced explicit supervision of well-stabilized sub-operations, some forms of flow may exhibit reduced activity in neural systems associated with self-monitoring or executive intervention, alongside preserved or increased activity in task-relevant control systems.
\end{hypothesis}

The qualification in Hypothesis~\ref{hyp:hypofrontality} is doing real work. Using perfusion and BOLD imaging during experimentally induced flow, Ulrich and colleagues found that flow was associated with \emph{decreased} activity in the medial prefrontal cortex and amygdala, but with \emph{increased} activity in the inferior frontal gyri, the anterior insula, the basal ganglia, and the midbrain \cite{ulrich2014,ulrich2016}---a pattern of selective reallocation among frontal regions rather than uniform prefrontal withdrawal. Contemporary reviews of the neuroscience of flow emphasize network-level reorganization, involving arousal and reward circuitry alongside frontal regions, over any single mechanism of blanket prefrontal deactivation \cite{vanderlinden2021}. The earlier draft of this essay stated the neural claim as a corollary that followed directly from the formal model; that was an overreach. The formal model in Section~\ref{subsec:art-formal} entails nothing about neural localization on its own; at most it motivates Hypothesis~\ref{hyp:hypofrontality} as a research question.

A related distinction from the earlier draft is worth keeping at the functional level rather than restating as an established neural equivalence. Reduced deliberative supervision can arise because supervision is no longer needed---a stabilized aspect continues to satisfy its constraints unattended, in the sense of Section~\ref{subsec:art-formal}---or because supervision is impaired while it is still needed, as under acute stress or threat arousal. The claim that these two conditions produce ``indistinguishable'' reductions in prefrontal activity is not supported: acute-stress effects on executive function vary with severity, timing, and task, and a meta-analysis finds that acute stress reliably impairs working memory and cognitive flexibility while producing more mixed, direction-dependent effects on inhibition \cite{shields2016}---a pattern that is not simply interchangeable with the frontal reallocation reported for flow. The functional distinction remains useful even without a claim of neural equivalence: supervision withdrawn because it is no longer needed, versus supervision impaired while it is still needed, is a difference in reliability and origin, not a difference that current neuroimaging demonstrates as a single shared signature.

\subsection{Hierarchical Attention and Compositional Depth}

Because attention is bounded, relegation does more than reduce felt effort. Under the revised update rule, freeing capacity in $A(t)$ via $R_t$ makes room for higher-order aspects to enter via $N_t$: coordination, abstraction, and integration across sub-tasks that were not representable while lower-level aspects still required monitoring. This produces a hierarchical attentional structure in which lower-level operations become largely automatic, mid-level structures are intermittently attended, and higher-level planning remains actively monitored, subject to $Q_t$ reactivating any level when context demands it. Skill, on this account, is not merely efficiency but depth of compositional control: the capacity to operate at a higher level of abstraction because lower levels have been sufficiently, though not irrevocably, stabilized.

\subsection{Skill as Compressed Structure}

\medskip
\begin{center}
\begin{tabular}{ll}
\toprule
\textbf{Structural Process} & \textbf{Cognitive Process} \\
\midrule
Expansion & Learning and exploration \\
Compression & Relegation ($R_t$) \\
Invariants & Stabilized aspects \\
Re-expansion & Reactivation ($Q_t$) \\
\bottomrule
\end{tabular}
\end{center}
\medskip

Relegated aspects encode structural invariants discovered through prior practice, and the felt ease of expert performance corresponds to a small active attention set relative to the full task, $|A(t)| \ll |T|$, rather than to any reduction in the task's underlying complexity. This has a direct implication for how ease should be read: that a task feels easy to an expert is evidence that the expert has performed enough exploration to stabilize its recurring structure, not evidence that the task is simple. The complexity has not disappeared; it has been internalized, and $Q_t$ is the model's reminder that this internalization is conditional, not permanent.

\subsection{Irreversibility and Reconstructive Loss}

Relegation introduces a partial loss of reconstructability. Once aspects have left $A(t)$, intermediate steps may no longer be explicitly representable, explanations may become compressed or incomplete, and reconstructing the full expansion history may require deliberate re-expansion. Expert knowledge is often operationally accessible but descriptively compressed: the expert can perform the task while struggling to narrate how. This is not a failure but a structural feature, and it explains why expert advice is so often correct for the expert and useless for the novice---the relegated aspects are no longer visible to the advisor, so the advice skips steps that are not yet stabilized for the listener. The compressed explanation presupposes an exploration the novice has not yet performed.

\subsection{An Extended Example: Expert and Novice Perception}
\label{subsec:expert-novice}

The divergence between expert and novice perception of a disordered system illustrates the model concretely. An expert surveying a derelict building perceives a tractable sequence of operations: assess structural integrity, identify constraint violations, decompose into sequenced sub-tasks, estimate materials and labor. A novice surveying the same building perceives disorder whose resolution is not apparent. Both perceptions are, in a precise sense, accurate about the building's current state; what differs is not information about $\Omega$ but the availability of a generative pathway toward $\Omega_{\mathrm{valid}}$. The expert's $K(t)$ already contains stabilized decompositions for this class of problem, so the perceived task is tractable rather than complete; the novice's $K(t)$ does not, and the system is, from the novice's own epistemic position, genuinely difficult to see a path through---not irrationally so, but appropriately so, given what has and has not yet been explored.

\subsection{An Extended Example: Writing as Expansion}
\label{subsec:writing}

Writing is a second useful illustration, because it makes the dependency argument of Section~\ref{subsec:dependency} concrete. A writer who attempts to produce the compressed form directly---the elegant, concise, finalized argument---without first performing the exploratory work will typically fail to produce it, not for want of talent but because compression under Definition~\ref{subsec:compression} presupposes having identified what is invariant. Drafting functions as expansion: articulating candidate structures, discovering dependencies, identifying failure modes in an argument. Revision functions as stabilization: correcting constraint violations, consolidating surviving structure, reducing redundancy. The final form is compression proper: encoding what expansion and stabilization have found. The common advice to ``write simply'' is not wrong as a description of a desirable endpoint, but it says little about the process by which that endpoint is ordinarily reached; for most writers, on most subjects, there is no shortcut through the exploratory phase, though a writer working in territory someone else has already mapped can sometimes begin closer to the compressed form by relying on that inherited exploration.

\section{Process Invisibility, Attribution, and Selection}
\label{sec:invisibility}

\subsection{Compressed Outputs and Shifted Burden}

The analysis of Sections~\ref{sec:formal}--\ref{sec:relegation} has a direct implication for how compressed outputs are interpreted in any domain. When a compressed representation is presented without its generative history, the interpretive burden shifts from producer to receiver: the producer has internalized the relevant exploration, and the receiver has not. This asymmetry has different consequences depending on context. In mathematics, an elegant proof may be genuinely illuminating for a reader who has worked in the domain, because the compression encodes structure that reader can reconstruct; the same proof may be opaque to a reader who lacks that background. The compression has not failed there; it has simply been addressed to a different receiver than the one encountering it.

In commercial and media contexts the situation is more troubling, because compressed outputs are frequently presented as though no prior exploration were required to reach them, when the audience's lack of background is precisely what the compression depends on to look effortless. The genre of compressed instruction---``just learn to code,'' ``just start writing,'' ``just build the product''---presents a relegated form as if it required no prior expansion, describing a valid terminal state while omitting the trajectory required to reach it.

There is a useful distinction between compression that is reconstructible by a prepared receiver and compression that is parasitic on undisclosed prior knowledge. Valid compression encodes structure an appropriately prepared receiver can, in principle, reconstruct; a minimal proof or a well-refactored codebase are examples, since the invariants they encode are recoverable given sufficient background. What might be called \emph{epistemic outsourcing} instead achieves an appearance of simplicity not by encoding structure efficiently but by silently delegating the exploratory work to a receiver who is not equipped to perform it and is not told they must. Where this pattern holds, observers can come to infer that production is as easy as consumption, feeding the attribution failures discussed next.

\subsection{Fundamental Attribution Error and the Erasure of Process}

This misperception is amplified by a well-established feature of social cognition. Ross's account of the fundamental attribution error describes a systematic tendency to attribute observed outcomes to dispositional properties of agents---talent, intelligence, character---while underweighting situational and processual factors \cite{ross1977}. Applied to skill and production, a successful outcome tends to be attributed to the producer's intrinsic qualities rather than to the structural conditions and accumulated exploration that made it possible; a failed attempt tends to be attributed to the producer's inadequacy rather than to genuine task difficulty. This pattern is not irrational given the available evidence---outcomes are visible, processes typically are not, and dispositional attribution is a reasonable response to a systematically incomplete information environment. But the environment is incomplete, and the resulting attributions inherit that incompleteness.

Media ecosystems can amplify this bias. Selection effects favor circulating outcomes that appear dramatically discontinuous with ordinary effort---the overnight success, the effortlessly productive artist---and narrative conventions tend to frame these as expressions of individual quality rather than structural process, while compression strips away the exploration history that would reveal the underlying structure. The sample of outcomes that reaches public attention is consequently not a representative sample of the underlying process, and treating it as such can produce miscalibrated expectations.

\subsection{Probabilistic Asymmetry and Survivor Narratives}

A related distortion operates at the level of population inference. In many domains of skilled production---entertainment, athletics, entrepreneurship, research, artistic practice---outcome distributions are heavily right-skewed: a small number of participants achieve highly visible success, and a large majority do not achieve comparable recognition regardless of comparable effort. The successes that circulate are consequently a non-representative sample of the underlying population of attempts, selected and amplified precisely because they are extreme. Canonization narratives constructed retrospectively around a small number of eventually-recognized figures tend to read as stories of difficult genius eventually rewarded; this reconstruction can suppress the contingency and selection effects that separated the canonized case from equally capable contemporaries whose recognition never arrived. This is offered as a structural pattern in how success narratives are typically assembled, not as a documented claim about any specific individual's biography, since establishing the latter would require evidence this essay does not undertake to marshal.

\subsection{Misattribution from Evaluative Ease}

The asymmetry formalized in Proposition~\ref{prop:asymmetry} and Section~\ref{subsec:fragility} has a direct bearing on this pattern of misattribution. Agents whose primary interaction with a system is evaluative encounter it already inside $\Omega_{\mathrm{valid}}$ and interact with it through the comparatively local task of locating a flaw. Where the apparent ease of that task is read as evidence that avoiding the flaw during construction was equally easy, the inference overreaches what Proposition~\ref{prop:asymmetry} actually supports---since, as noted in Section~\ref{subsec:asymmetry}, detecting or explaining a failure is not always cheap, and reviewing a proof or diagnosing a distributed system can be more demanding than producing a locally plausible artifact. What the formal result does support is narrower: construction and evaluation typically sample from different, and often differently sized, regions of the space of possible operations, so evaluative ease is not on its own strong evidence about constructive difficulty, in either direction.

\section{Labor and the Political Economy of Delegated Constraint Management}
\label{sec:labor}

\subsection{The Visibility Gap}

The preceding analysis converges on an asymmetry in how different categories of labor tend to be perceived and rewarded. Essential work---infrastructural, repeatable, stabilizing labor such as maintenance, care, teaching, and construction---often manifests as the absence of failure rather than the presence of a discrete success, and can therefore go relatively unnoticed even when it sustains the conditions under which other, more visible production becomes possible. Visible work, by contrast, tends to be rewarded roughly in proportion to its visibility, benefiting from the amplification mechanisms described in Section~\ref{sec:invisibility}. Recognition of essential work is also often temporally displaced, arriving---if it arrives at all---after some later, more visible event makes the earlier contribution retrospectively legible; a teacher may be recognized only through the lens of a student's later prominence. This displacement can feed back into the attribution pattern of Section~\ref{sec:invisibility}, biasing the cultural model of achievement toward visible compression and away from the labor that preceded it.

\subsection{Delegation, Constraint Blindness, and Compressed Advice}

Wealth confers, among other things, the capacity to delegate constraint management: supply chains, maintenance schedules, coordination problems that structure everyday life can be handled by others and rendered invisible to the person delegating them. This delegation is not inherently problematic, and its epistemic consequence is a tendency rather than a certainty: delegation \emph{can} produce a kind of constraint blindness, a reduced capacity to perceive the structural conditions that make an apparent solution non-trivial, though it need not, and the degree to which it does varies with how the delegating agent remains engaged with the delegated domain. Where it does occur, it tends to produce a familiar genre of compressed advice---``just do this,'' ``simply approach it differently''---that is not always dishonest so much as reflective of a genuine perceptual gap between a speaker whose constraints are managed and a listener whose are not. Such advice can be harmful when treated as actionable guidance, since it presents a compressed form that presupposes an exploration the listener must still perform but has not been told about.

\section{The Ethics of Compression}
\label{sec:ethics}

\subsection{Three Distinct Standards}

The arguments developed above ground a normative point about the ethics of representation, but the earlier draft's standard---that an ethically compliant compressed representation must preserve, in principle, the reconstructibility of its actual generative pathway---is too demanding and in some cases incoherent. A finished mathematical proof need not preserve the psychological history of its discovery; a reader who independently reconstructs a valid derivation has not thereby recovered the author's actual genesis, and does not need to. Three distinct standards should be kept apart.

\emph{Logical reconstructibility} asks whether a prepared receiver could, in principle, derive or verify the compressed content from stated premises, independent of how it was actually produced. \emph{Practical teachability} asks whether the compressed form, or some accompanying material, gives an unprepared receiver a workable path toward the competence the compression represents. \emph{Historical provenance} asks whether the compressed form accurately represents the actual process---whose labor, how much exploration, over what period---that produced it.

\subsection{Non-Deception as the Working Standard}

Requiring all three uniformly is too strong: mathematics regularly and legitimately satisfies logical reconstructibility while providing neither teachability nor provenance, and this is not an ethical defect in a proof. What the arguments of this essay actually support is narrower: a compressed representation becomes ethically problematic when it is used to \emph{misrepresent} the labor or prerequisites behind it---when a producer or distributor actively implies that the compressed form required no meaningful exploration, when it obscures whose labor performed that exploration, or when it is marketed as teachable when it is not.

\begin{definition}
A compressed representation is \emph{non-deceptive} if it does not affirmatively misrepresent the exploration, labor, or prerequisites that produced it, whether or not it also satisfies logical reconstructibility or practical teachability.
\end{definition}

\subsection{Criteria for Ethical Representation}

Three working criteria follow from this narrower standard. First, non-deception about labor: representations of skilled output should not imply, where it is false, that no meaningful prior exploration was required. Second, appropriate matching of standard to context: a representation should be evaluated against the standard---logical, pedagogical, or historical---that it actually claims to meet, not against all three at once. Third, accuracy of effort distribution: representations should not systematically imply that the observed outcome is more typical of the underlying attempt distribution than it is. These criteria are context-dependent rather than absolute, but they ground a substantive critique of representational practices in marketing, pedagogy, and media that violate them, without requiring the stronger and less defensible claim that every compressed artifact owes its receiver a full account of its own history.

\section{Conclusion: Toward a Narrower and More Defensible Claim}
\label{sec:conclusion}

The argument of this essay can be summarized in a form more modest than its opening thesis, and correspondingly harder to refute. Under a stated sparse-validity assumption and a stated assumption about perturbation distributions, maintaining coherence in a system requires selecting from a comparatively narrow set of operations, while an unstructured perturbation is likely to leave that set by an easier route; this is a conditional probabilistic result, not a universal entropy law. In human skill, a substantial and well-evidenced class of expert automaticity arises from procedures that were once deliberately controlled and became increasingly autonomous through practice, chunking, and probabilistic stabilization, though this class does not exhaust automatic cognition, much of which has other origins entirely. Observers who encounter only a compressed result, without its generative history, are consequently prone---not certain, but prone---to underestimate the exploration that produced it, and this misattribution has downstream consequences for how labor is valued, how advice is given, and how expectations are formed.

The correction proposed here is not a recommendation to make things more complicated, nor a claim that difficulty is always hidden or that ease is always suspect. It is a recommendation for epistemic caution about the structure of generative processes: an acknowledgment that a compressed artifact, presented without its exploration history, is frequently a partial and potentially misleading representation of the work that produced it, conditional on the assumptions this essay has tried to state rather than assume. Understanding this structure does not make exploration easier. It does make its erasure a little harder to take at face value.

\newpage
\appendix

\section{Summary of Formal Results}

\begin{description}
\item[Proposition~\ref{prop:compression-requires-expansion}] (Conditional, epistemic.) In an unknown domain, a reliable compressed representation generally requires prior exploration by some agent, though not necessarily the one presenting the compression.
\item[Proposition~\ref{prop:asymmetry}] (Conditional on Assumption~\ref{ass:sparse} and a stated perturbation-distribution assumption.) A random perturbation of a valid configuration is more likely to leave the valid region than to remain within it. This does not by itself establish that evaluation is easier than construction.
\item[Proposition~\ref{prop:s1-as-s2}] (Restricted scope.) A substantial class, not the entirety, of expert Type 1 performance derives from relegated, formerly deliberate procedures.
\item[Hypothesis~\ref{hyp:hypofrontality}] (Contested, not established.) Some forms of flow may involve reduced activity in self-monitoring systems alongside preserved or increased task-relevant control activity; the empirical picture is one of selective reallocation, not uniform prefrontal withdrawal, and is actively debated.
\end{description}

\section{Glossary of Key Terms}

\begin{description}
\item[Aspect] A constraint, sub-operation, or dependency constituting part of a task.
\item[Active attention set $A(t)$] The subset of known aspects under conscious monitoring at time $t$, bounded by capacity $k$.
\item[Known aspects $K(t)$] The subset of a task's aspects an agent has represented, whether or not currently attended.
\item[Coherent configuration] A configuration satisfying all constraints in $\mathcal{C}$.
\item[Compression $(C, D, \sim)$] An encoder, decoder, and task-relevant equivalence relation such that decoding an encoding recovers the original up to $\sim$.
\item[Sparse validity] The stated assumption that valid configurations occupy small measure within the full configuration space, for the systems this essay discusses.
\item[Expansion] A process of configuration-space traversal through which invariant structure is discovered or transmitted.
\item[Relegation ($R_t$)] The subset of the active attention set withdrawn at time $t$ because it has met a stabilization threshold.
\item[Reactivation ($Q_t$)] Aspects, possibly previously relegated, returned to active attention by novelty, error, or unfamiliar context.
\item[Stabilization] A probabilistic, context-relative condition: an aspect satisfies its constraint with probability at least $1-\varepsilon$ under a given policy and environment class.
\end{description}

\begin{thebibliography}{99}

\bibitem{bargh1992} Bargh, J. A. (1992). The ecology of automaticity: Toward establishing the conditions needed to produce automatic processing effects. \textit{American Journal of Psychology}, 105(2), 181--199.

\bibitem{bargh2012} Bargh, J. A., Schwader, K. L., Hailey, S. E., Dyer, R. L., \& Boothby, E. J. (2012). Automaticity in social-cognitive processes. \textit{Trends in Cognitive Sciences}, 16(12), 593--605.

\bibitem{dietrich2003} Dietrich, A. (2003). Functional neuroanatomy of altered states of consciousness: The transient hypofrontality hypothesis. \textit{Consciousness and Cognition}, 12(2), 231--256.

\bibitem{dietrich2004} Dietrich, A. (2004). Neurocognitive mechanisms underlying the experience of flow. \textit{Consciousness and Cognition}, 13(4), 746--761.

\bibitem{evans2013} Evans, J. St. B. T., \& Stanovich, K. E. (2013). Dual-process theories of higher cognition: Advancing the debate. \textit{Perspectives on Psychological Science}, 8(3), 223--241.

\bibitem{kal2018} Kal, E., Prosée, R., Winters, M., \& van der Kamp, J. (2018). Does implicit motor learning lead to greater automatization of motor skills compared to explicit motor learning? A systematic review. \textit{PLOS ONE}, 13(9), e0203591.

\bibitem{kahneman2011} Kahneman, D. (2011). \textit{Thinking, Fast and Slow}. Farrar, Straus and Giroux.

\bibitem{ross1977} Ross, L. (1977). The intuitive psychologist and his shortcomings: Distortions in the attribution process. In L. Berkowitz (Ed.), \textit{Advances in Experimental Social Psychology} (Vol. 10, pp. 173--220). Academic Press.

\bibitem{shields2016} Shields, G. S., Sazma, M. A., \& Yonelinas, A. P. (2016). The effects of acute stress on core executive functions: A meta-analysis and comparison with cortisol. \textit{Neuroscience \& Biobehavioral Reviews}, 68, 651--668.

\bibitem{ulrich2014} Ulrich, M., Keller, J., Hoenig, K., Waller, C., \& Grön, G. (2014). Neural correlates of experimentally induced flow experiences. \textit{NeuroImage}, 86, 194--202.

\bibitem{ulrich2016} Ulrich, M., Keller, J., \& Grön, G. (2016). Neural signatures of experimentally induced flow experiences identified in a typical fMRI block design with BOLD imaging. \textit{Social Cognitive and Affective Neuroscience}, 11(3), 496--507.

\bibitem{vanderlinden2021} Van der Linden, D., Tops, M., \& Bakker, A. B. (2021). Go with the flow: A neuroscientific view on being fully engaged. \textit{European Journal of Neuroscience}, 53(4), 947--963.

\end{thebibliography}

\end{document}

